\documentclass{article}
\usepackage{amsmath}
\usepackage{amsthm}
\usepackage{siunitx}
\usepackage{booktabs}
\usepackage[margin=1in]{geometry}
\newtheorem{remark}{Remark}
\sisetup{detect-all}

\title{R-32 Vapor-Compression Cycle on the IDAES Peng--Robinson Package:\\
State Definition, Degrees of Freedom, Phase Selection and Results}
\author{Shilpa Narasimhan}
\date{October 15, 2026}

\begin{document}
\maketitle

\section{Scope}
This note documents how the single-stage R-32 cycle is set up and computed on
the IDAES generic Peng--Robinson (PR) property package with the \texttt{FPhx}
state definition, how its states are seeded and checked, and the results to
date. The PR equation of state, the departure functions and the Shomate
ideal-gas terms are as in \texttt{R32\_DVCT\_computations\_writeup.tex} and are
not repeated here. Code: \texttt{vapor\_compression\_cubic\_FPhx\_1015.py},
\texttt{closedform\_4methods\_1015.py}, \texttt{probe\_vle\_1015.py},
\texttt{probe\_bound\_1015.py}, \texttt{run\_FPhx\_cycle\_1015.py},
\texttt{sweep\_FPhx\_1015.py}.

\section{Cycle Specification}
\begin{table}[h]
\centering
\begin{tabular}{ll}
\toprule
Quantity & Value \\
\midrule
Evaporating temperature $T_e$ & \SI{-29}{\celsius} \\
Condensing temperature $T_c$ & $T_{amb} + \SI{9}{\kelvin}$ \\
Ambient temperatures $T_{amb}$ & \SIlist{10;15;20;25}{\celsius} \\
Superheat $\Delta T_{SH}$ & \SI{3}{\kelvin} (window \SIrange{3}{8}{\kelvin}) \\
Subcooling $\Delta T_{SC}$ & \SI{8}{\kelvin} (window \SIrange{3}{8}{\kelvin}) \\
Compressor isentropic efficiency $\eta$ & 0.80 \\
Refrigerant mass flow & \SI{1}{\kilo\gram\per\second} \\
Expansion valve & isenthalpic \\
\bottomrule
\end{tabular}
\caption{Operating specification (DVCT).}
\end{table}

\section{Property Methods}
Three PR parameter sets are compared with the Helmholtz reference
(CoolProp, Tillner-Roth and Yokozeki). The ideal-gas heat capacity is in Shomate form,
$C_p^{ig} = A + Bt + Ct^2 + Dt^3 + E/t^2$, $t = T/\SI{1000}{\kelvin}$.

\begin{table}[h]
\centering
\begin{tabular}{lrrrrrrrr}
\toprule
Method & $P_c$ [bar] & $T_c$ [K] & $\omega$ & $A$ & $B$ & $C$ & $D$ & $E$ \\
\midrule
NIST & 57.82 & 351.300 & 0.2769 & $-6.0987$ & 179.220 & $-122.368$ & 32.302 & 0.4914 \\
GCGP (scaled) & 50.73 & 355.354 & 0.1711 & 14.161 & 124.0 & $-63.40$ & 11.90 & 0 \\
SPGP (\texttt{fp\_single\_0919}) & 52.168 & 325.387 & 0.3523 & 32.6437 & 26.9797 & 591.354 & $-597.447$ & $-0.04746$ \\
\bottomrule
\end{tabular}
\caption{PR parameters. GCGP coefficients converted to Shomate: $A=a$,
$B=b\times10^3$, $C=c\times10^6$, $D=d\times10^9$, $E=e\times10^{-6}$.
SPGP $\omega$ from $P^{sat}=\SI{1738.797}{mmHg}$ at $T_r=0.7$:
$\omega = -1-\log_{10}(P^{sat}/P_c)$. IDAES uses $R=\SI{8.314462618}{\joule\per\mole\per\kelvin}$.}
\end{table}

\section{Phase-Specific Properties}
At a given $(T,P)$ the PR cubic has either one admissible real root or three
real roots. IDAES evaluates every property separately for each phase $p\in\{L,V\}$:
the liquid phase takes the smallest root $Z_L$, the vapor phase the largest
root $Z_V$. Outside the three-root region both phases receive the same root,
so $h_L = h_V$ and $\phi_L = \phi_V$. For each phase,
\begin{equation}
h_p(T,P) = h^{ig}(T) + h^{dep}(T,P,Z_p), \qquad
s_p(T,P) = s^{ig}(T,P) + s^{dep}(T,P,Z_p).
\end{equation}
The mixture properties of a state block are phase-fraction weighted,
\begin{equation}\label{eq:mix}
h = \beta_V\, h_V(T,P) + (1-\beta_V)\, h_L(T,P),
\end{equation}
where $\beta_V$ is the vapor fraction (\texttt{phase\_frac["Vap"]}).

\subsection{Choosing the Physical Phase Without a Solver}
For a pure component at fixed $(T,P)$ the stable phase is the one with the
lower fugacity. Define
\begin{equation}
\Delta(T,P) = \ln\phi_L(T,P) - \ln\phi_V(T,P).
\end{equation}
$\Delta < 0$: liquid is stable. $\Delta > 0$: vapor is stable. $\Delta = 0$ with
three roots: saturation. In code (\texttt{PRProbe}), $h_p$ and $s_p$ are read directly
from the IDAES expressions \texttt{enth\_mol\_phase[p]} and
\texttt{entr\_mol\_phase[p]} for the phase chosen by the sign of $\Delta$.
No equation is solved, so a wrong root cannot be returned.

\subsection{PR's Own Saturation Line}
The saturation pressure and temperature are those of the PR model itself,
\begin{equation}
\Delta\big(T, P^{sat}_{PR}(T)\big) = 0,
\end{equation}
found by bisection. Outside the three-root window $\Delta$ is exactly zero, so
the bracket is first located by scanning for a change of sign between
\emph{nonzero} values of $\Delta$; a plain bisection drifted to $P_c$ for SPGP.

\begin{table}[h]
\centering
\begin{tabular}{lrrrr}
\toprule
& \multicolumn{4}{c}{$P^{sat}$ [kPa]} \\
$T$ & Helmholtz & NIST & GCGP & SPGP \\
\midrule
\SI{-29}{\celsius} & 284.9 & 283.3 & 287.8 & 469.7 \\
\SI{19}{\celsius} & 1434.2 & 1443.1 & 1291.3 & 2289.3 \\
\SI{24}{\celsius} & 1644.8 & 1656.4 & 1467.7 & 2619.9 \\
\SI{29}{\celsius} & 1878.0 & 1892.6 & 1661.2 & 2985.4 \\
\SI{34}{\celsius} & 2135.3 & 2153.1 & 1873.1 & 3388.2 \\
\bottomrule
\end{tabular}
\caption{Saturation pressures at the evaporating and condensing temperatures.}
\end{table}

\begin{remark}
Because each method has its own saturation line, a state that is vapor in one
method can be liquid in another. Helmholtz states therefore cannot be used to
seed PR states directly: $T$ and $P$ transfer, but $h$ does not (different
reference state) and the phase does not (e.g.\ SPGP's condensing pressure is
\SI{60}{\percent} above Helmholtz's).
\end{remark}

\section{Closed-Form Cycle}\label{sec:closed}
The reference COP for each PR method uses the same PR expressions as the
IDAES model, but no flowsheet and no simultaneous solve.
\begin{align}
P_e &= P^{sat}_{PR}(T_e), \qquad P_c = P^{sat}_{PR}(T_{cond})\\
\text{State 1 (evaporator outlet):}\quad
T_1 &= T_e + \Delta T_{SH}, \quad h_1 = h_V(T_1,P_e),\quad s_1 = s_V(T_1,P_e)\\
\text{State 2s (isentropic):}\quad
&s_V(T_{2s},P_c) = s_1 \;\Rightarrow\; h_{2s} = h_V(T_{2s},P_c)\\
\text{State 2 (compressor outlet):}\quad
h_2 &= h_1 + \frac{h_{2s}-h_1}{\eta}\\
\text{State 3 (condenser outlet):}\quad
T_3 &= T_{cond} - \Delta T_{SC}, \quad h_3 = h_L(T_3,P_c)\\
\text{State 4 (valve outlet):}\quad
h_4 &= h_3\\
\mathrm{COP} &= \frac{h_1 - h_4}{h_2 - h_1}
\end{align}
$T_{2s}$ and $T_2$ are found by bisection on the vapor branch. The Helmholtz
COP uses the same steps with CoolProp. The Carnot limit is
$\mathrm{COP}_{Carnot} = T_e/(T_{cond}-T_e)$ in kelvin.

\begin{remark}
This reference checks the flowsheet and the solver, not the PR parameters:
both use the same PR equations. The check on the parameters is the
comparison with Helmholtz.
\end{remark}

\begin{remark}
The earlier closed form (\texttt{cop\_vs\_ambient\_closedform\_1014.py}) is
not pure PR: its saturation pressure comes from a corresponding-states
correlation, $\omega$ from interpolated Linde data, and $P_c = \SI{57.85}{bar}$.
At $T_{amb}=\SI{20}{\celsius}$ it gives 2.6243 for NIST against 2.6149 here
(\SI{0.36}{\percent}).
\end{remark}

\section{IDAES Flowsheet}
\subsection{Units and State Definition}
Four units in a closed loop: evaporator (\texttt{Heater}), compressor
(\texttt{Compressor}, isentropic), condenser (\texttt{Heater}), expansion valve
(\texttt{PressureChanger}, adiabatic, $\Delta P<0$). The property package is the
generic PR package with the \texttt{FPhx} state definition:
\begin{equation}
\text{state variables} = \{F,\; x,\; P,\; h\}.
\end{equation}
Each arc therefore writes equalities on $F$, $x$, $P$ and $h$. Because
$(P,h)$ determines the thermodynamic state of a pure fluid everywhere,
including inside the dome, the stream leaving one unit is the stream entering
the next.

\begin{remark}
Under \texttt{FTPx} the arcs carry $(F,x,T,P)$. For a pure fluid at saturation
$T$ and $P$ are not independent, so the vapor fraction was not transmitted and
the loop did not conserve energy (breadcrumb, 10/08). The Helmholtz model
avoids this with \texttt{StateVars.TPX}, whose ``X'' is vapor fraction. In
the generic package, ``x'' in \texttt{FTPx} is mole fraction, which is always
1 for a pure fluid.
\end{remark}

\subsection{Degrees of Freedom}
DOF is read from the built model with \texttt{degrees\_of\_freedom()}:
\begin{table}[h]
\centering
\begin{tabular}{lr}
\toprule
Model state & DOF \\
\midrule
Units and arcs, nothing fixed & 5 \\
Fix $F$, $x$, $P_e$, $P_c$, $\eta$, $T_1$, $T_3$ & $-2$ \\
Deactivate one arc's $F$ and $x$ equalities (closed loop) & 0 \\
Release $T_1$, $T_3$ (superheat, subcooling) & \textbf{2} \\
\bottomrule
\end{tabular}
\end{table}
In a closed loop every unit conserves flow and composition, so one arc's
flow and composition equalities are implied by the other three; exactly those
two are removed. Mass continuity is verified afterwards by
\texttt{cycle\_guards}. The two remaining degrees of freedom are the
evaporator and condenser outlet temperatures:
\begin{equation}
T_1 = T^{sat}_{PR}(P_e) + \Delta T_{SH}, \qquad
T_3 = T^{sat}_{PR}(P_c) - \Delta T_{SC}.
\end{equation}

\section{The Pure-Component Phase Problem}
\subsection{Observation}
A single state block was solved at known states with $P$ and $h$ (or $T$ and
$P$) fixed. Every combination of VLE formulation (\texttt{SmoothVLE},
\texttt{CubicComplementarityVLE}) and state definition (\texttt{FTPx},
\texttt{FPhx}) returned impossible states and reported them optimal, for
example:
\begin{table}[h]
\centering
\begin{tabular}{llllrrr}
\toprule
State & VLE & Def. & Start & $T$ [\si{\celsius}] & True $T$ & $\beta_V$ \\
\midrule
Vapor, $P=\SI{283.3}{kPa}$ & Smooth & FPhx & at truth & 50.43 & $-26.00$ & 0.729 \\
Liquid, $P=\SI{1878}{kPa}$ & Smooth & FPhx & at truth & 17.30 & 21.00 & 0.027 \\
Liquid, $P=\SI{1878}{kPa}$ & Smooth & FTPx & default & 21.00 & 21.00 & 0.278 \\
Liquid, $P=\SI{1878}{kPa}$ & CubicCC & FPhx & default & 17.92 & 21.00 & 0.023 \\
\bottomrule
\end{tabular}
\caption{Isolated state blocks, NIST. Started at the true state, the solver
moves away from it.}
\end{table}

\subsection{Explanation}
For a mixture, the liquid and vapor phases have different compositions, and
the component balances fix $\beta_V$. For a pure component every composition
is 1, so the component balances give no information about $\beta_V$. Away
from saturation, $h_L \neq h_V$ at a given $T$ (three roots), and
Eq.~\eqref{eq:mix} can be satisfied by a blend of the two roots at a
\emph{wrong} temperature. The block's equations are square by count, but this
family of solutions is not excluded. Physically, away from saturation
$\beta_V$ must be exactly 0 or 1, and inside the dome $T$ must equal
$T^{sat}$.

\section{Remedy: Bounds}
Bounds do not change the degrees of freedom.

\subsection{Single-Phase Blocks}
\begin{equation}
\beta_V \ge 1-\varepsilon \;\;\text{(vapor blocks)}, \qquad
\beta_V \le \varepsilon \;\;\text{(liquid blocks)}.
\end{equation}
Vapor: evaporator outlet, compressor inlet and outlet, condenser inlet,
compressor isentropic block. Liquid: condenser outlet, valve inlet.
$\varepsilon$ was chosen by measurement (\texttt{probe\_bound\_1015.py}, NIST,
\texttt{SmoothVLE}, \texttt{FPhx}):
\begin{table}[h]
\centering
\begin{tabular}{lrrrr}
\toprule
& \multicolumn{4}{c}{$|T - T_{true}|$ [K]} \\
$\varepsilon$ & Vapor \SI{-26}{\celsius} & Liquid \SI{21}{\celsius} & Liquid \SI{26}{\celsius} & Liquid \SI{14}{\celsius} \\
\midrule
$10^{-3}$ & 0.320 & 0.089 & (infeasible) & 0.096 \\
$10^{-4}$ & 0.034 & 0.009 & 0.008 & 0.010 \\
$\mathbf{10^{-5}}$ & \textbf{0.004} & \textbf{0.001} & \textbf{0.001} & \textbf{0.001} \\
$10^{-6}$ & 0.0003 & (infeasible) & (infeasible) & 0.0002 \\
\bottomrule
\end{tabular}
\caption{The temperature error falls tenfold for each tenfold tightening of
$\varepsilon$. $\varepsilon = 10^{-5}$ is the tightest value that solved
every case. \texttt{CubicComplementarityVLE} with bounds was infeasible in
every case.}
\end{table}

\subsection{Two-Phase Blocks}
The valve outlet and evaporator inlet are two-phase. Their pressure is fixed
($P_e$), so $T^{sat}_{PR}(P_e)$ is known in advance:
\begin{equation}
T^{sat}_{PR}(P_e) - \delta \le T \le T^{sat}_{PR}(P_e) + \delta, \qquad \delta = \SI{0.01}{\kelvin}.
\end{equation}
\begin{remark}
$\delta$ has not been tested the way $\varepsilon$ was, and it is binding:
the solved temperature sits at the upper edge, \SI{-28.99}{\celsius}.
\end{remark}

\section{Seeding}
\begin{enumerate}
\item Compute the closed-form states (Section~\ref{sec:closed}).
\item Set $F$, $x$, $P$, $h$, $T$ and $\beta_V$ of every state block from those
states, with the phase chosen by the sign of $\Delta$. Two-phase blocks get
$T = T^{sat}_{PR}(P_e)$ and the lever-rule quality
$\beta_V = (h_4 - h_L^{sat})/(h_V^{sat}-h_L^{sat})$.
\item Apply the $\beta_V$ bounds.
\item Solve each state block alone at fixed $(F,x,P,h)$.
\item Apply the two-phase temperature bounds and solve the full flowsheet.
\end{enumerate}
Outlet state blocks carry a redundant $\sum x = 1$ equation, so step 4 is
refused (``too few degrees of freedom'') on those blocks, which keep their
step-2 values. IDAES's own state-block \texttt{initialize()} is not used: it
ignores the bounds.

\section{Checks Applied to Every Solved Cycle}
\begin{enumerate}
\item Carnot: $\mathrm{COP} \le T_e/(T_{cond}-T_e)$.
\item Loop energy closure: $Q_{evap} + W_{comp} + Q_{cond} = 0$.
\item Arc continuity of $h$ and $F$.
\item Phase check on all eight state blocks: $\beta_V$ on the side of
$T^{sat}_{PR}(P)$ that $T$ is on; two-phase allowed within \SI{0.05}{\kelvin} of it.
\item COP against the closed form from the same PR equations.
\end{enumerate}

\section{Results}
\subsection{IDAES PR Cycle, NIST, $T_{amb}=\SI{20}{\celsius}$}
\begin{table}[h]
\centering
\begin{tabular}{lrrrr}
\toprule
Block & $T$ [\si{\celsius}] & $T^{sat}_{PR}$ [\si{\celsius}] & $\beta_V$ & $h$ [\si{\joule\per\mole}] \\
\midrule
Evaporator inlet & $-28.99$ & $-29.00$ & 0.2525 & $-13338.6$ \\
Evaporator outlet & $-26.00$ & $-29.00$ & 1.0000 & 1084.9 \\
Compressor outlet & 114.30 & 29.00 & 1.0000 & 6600.8 \\
Condenser outlet & 21.00 & 29.00 & 0.0000 & $-13338.6$ \\
Valve outlet & $-28.99$ & $-29.00$ & 0.2525 & $-13338.6$ \\
\bottomrule
\end{tabular}
\caption{DOF 0, IPOPT optimal. COP 2.6149, closed form 2.6149
($\SI{-0.0001}{\percent}$), Carnot 4.2095, loop closure \SI{0.01}{W}. All
checks pass. Reproduced on DOWLINGLAB-7.}
\end{table}

\subsection{Closed-Form COP, All Methods}
\begin{table}[h]
\centering
\begin{tabular}{rrrrrr}
\toprule
$T_{amb}$ [\si{\celsius}] & Carnot & Helmholtz & NIST & GCGP & SPGP \\
\midrule
10 & 5.0865 & 3.3789 & 3.3540 & 3.4116 & 3.2058 \\
15 & 4.6066 & 2.9805 & 2.9515 & 3.0063 & 2.7615 \\
20 & 4.2095 & 2.6484 & 2.6149 & 2.6671 & 2.3818 \\
25 & 3.8754 & 2.3667 & 2.3282 & 2.3779 & 2.0499 \\
\midrule
\multicolumn{2}{l}{vs.\ Helmholtz} & --- & $-0.7$ to $-1.6\,\%$ & $+0.5$ to $+1.0\,\%$ & $-5$ to $-13\,\%$ \\
\bottomrule
\end{tabular}
\caption{$\Delta T_{SH}=\SI{3}{K}$, $\Delta T_{SC}=\SI{8}{K}$, $\eta=0.80$.
All below Carnot; all decrease monotonically with ambient. The Helmholtz
column matches the IDAES Helmholtz model maximised over superheat and
subcooling (\texttt{sweep\_helmholtz\_maximized\_1015.csv}), which lands at
$\Delta T_{SH}=\SI{3}{K}$, $\Delta T_{SC}=\SI{8}{K}$.}
\end{table}

\begin{table}[h]
\centering
\begin{tabular}{rrrrr}
\toprule
$T_{amb}$ [\si{\celsius}] & Helmholtz & NIST & GCGP & SPGP \\
\midrule
10 & 85.5 & 92.7 & 76.8 & 44.7 \\
15 & 96.0 & 103.6 & 86.3 & 51.3 \\
20 & 106.3 & 114.3 & 95.6 & 57.9 \\
25 & 116.4 & 124.8 & 104.7 & 64.4 \\
\bottomrule
\end{tabular}
\caption{Compressor discharge temperature [\si{\celsius}], closed form.}
\end{table}

\subsection{IDAES PR Sweep, Superheat and Subcooling Fixed}
\begin{table}[h]
\centering
\begin{tabular}{llrrrrcc}
\toprule
Method & $T_{amb}$ & IPOPT & COP & Closed form & Diff [\%] & Phases & Guards \\
\midrule
NIST & 10 & infeasible & 3.3540 & 3.3540 & $-0.0003$ & ok & ok \\
NIST & 15 & infeasible & 2.9515 & 2.9515 & $-0.0003$ & ok & ok \\
NIST & 20 & optimal & 2.6149 & 2.6149 & $-0.0001$ & ok & ok \\
NIST & 25 & optimal & 2.3282 & 2.3282 & $-0.0002$ & ok & ok \\
GCGP & 10 & optimal & 3.4116 & 3.4116 & $-0.0002$ & ok & ok \\
GCGP & 15 & optimal & 3.0063 & 3.0063 & $-0.0003$ & ok & ok \\
GCGP & 20 & optimal & 2.6671 & 2.6671 & $-0.0003$ & ok & ok \\
GCGP & 25 & optimal & 2.3779 & 2.3779 & $-0.0003$ & ok & ok \\
SPGP & 10 & infeasible & 3.2057 & 3.2058 & $-0.0007$ & ok & ok \\
SPGP & 15 & optimal & 2.7615 & 2.7615 & $-0.0007$ & ok & ok \\
SPGP & 20 & infeasible & 2.3818 & 2.3818 & $-0.0008$ & ok & ok \\
SPGP & 25 & infeasible & 2.0499 & 2.0499 & $-0.0008$ & ok & ok \\
\bottomrule
\end{tabular}
\caption{\texttt{sweep\_FPhx\_1015.py square}, DOWLINGLAB-7, \texttt{myidaesenv}.
$\Delta T_{SH}=\SI{3}{K}$, $\Delta T_{SC}=\SI{8}{K}$, DOF 0, loop closure
$\le \SI{0.02}{W}$ in every case. Five of twelve points end ``infeasible''
while matching the closed form to within \SI{0.001}{\percent} and passing every check.}
\end{table}

\section{Open Items}
\begin{enumerate}
\item Five ``infeasible'' points in the square sweep: the unsatisfied
equations and their residuals have not yet been examined. Suspect: a bound
fighting an equation (the two-phase $\delta$ is already binding at
\SI{20}{\celsius}).
\item Sweep with superheat and subcooling free (\texttt{sweep\_FPhx\_1015.py optimize}).
\item Test the two-phase temperature bound $\delta$ as $\varepsilon$ was tested.
\item Seeding: deactivate the redundant $\sum x=1$ on outlet blocks during
step 4, and apply the two-phase bound before step 4.
\end{enumerate}

\end{document}
